% الجزء: sets-functions-relations
% الفصل: sets
% القسم: russells-paradox

\documentclass[../../../include/open-logic-section]{subfiles}


\begin{document}

\olfileid[ar]{sfr}{set}{rus}
\olsection{مفارقة راسل}

تسوّغ الامتدادية الترميز $\Setabs{x}{\phi(x)}$ للدلالة على
\emph{مجموعة} قيم $x$ التي تحقق~$\phi(x)$. غير أن كل ما تسوّغه
الامتدادية \emph{حقًا} هو الفكرة الآتية. \emph{إذا} وُجدت مجموعة
عناصرها هي جميع الأشياء التي تحقق~$\phi$ ولا شيء سواها، \emph{فإنها}
المجموعة الوحيدة من هذا القبيل. وبعبارة أخرى، بعد تثبيت خاصية ما~$\phi$،
تكون المجموعة $\Setabs{x}{\phi(x)}$ فريدة، \emph{إن وُجدت}.

لكن هذا الشرط مهم!{} والأهم أن بعض الخواص لا تقبل
\emph{الاستيعاب}. أي إن بعض الخواص \emph{لا} تعرّف مجموعات.
ولو كانت جميعها تعرّف مجموعات، لوقعنا في تناقضات صريحة. وأشهر مثال
على ذلك مفارقة راسل.

قد تكون المجموعات نفسها !!{element}s في مجموعات أخرى---فمجموعة القوى
لمجموعة~$A$، مثلًا، مؤلفة من مجموعات. ولذلك من المعقول أن نسأل أو
نبحث هل مجموعة ما هي !!a{element} من مجموعة أخرى. فهل يمكن أن تكون
مجموعة ما عضوًا في نفسها؟ لا يبدو أن في فكرة المجموعة ما يستبعد ذلك.
فمثلًا، إذا كانت \emph{جميع} المجموعات تؤلف تجمّعًا من الأشياء، فقد
يُظن أنه يمكن جمعها في مجموعة واحدة---هي مجموعة جميع المجموعات.
وهذه المجموعة، لكونها مجموعة، لكانت !!a{element} من مجموعة جميع
المجموعات.

تنشأ مفارقة راسل عندما ننظر في خاصية ألا تكون المجموعة !!a{element}
من نفسها، أي أن تكون \emph{غير منتمية إلى نفسها}. فماذا لو
افترضنا وجود مجموعة جميع المجموعات التي لا تكون كل واحدة منها
!!a{element} من نفسها؟ فهل المجموعة
\[
R = \Setabs{x}{x \notin x}
\]
موجودة؟ يتبين أننا نستطيع إثبات عدم وجودها.

\begin{thm}[مفارقة راسل]\ollabel{thm:russells-paradox}
	لا توجد مجموعة $R = \Setabs{x}{x \notin x}$.
\end{thm}

\begin{proof}
إذا وُجدت $R = \Setabs{x}{x \notin x}$، فإن
$R \in R$ إذا وفقط إذا كان $R \notin R$، وهذا تناقض.
\end{proof}

\begin{tagblock}{novice}
\begin{explain}
لنتناول هذا البرهان بتمهّل أكبر. إذا وُجدت~$R$، كان من المعقول أن نسأل
هل $R \in R$ أم لا. افترض أن $R \in R$ بالفعل. لقد عُرّفت~$R$ بأنها
مجموعة جميع المجموعات التي ليست !!{element}s في نفسها. فإذا كان
$R \in R$، لم تكن~$R$ نفسها متصفة بخاصية~$R$ التعريفية. لكن المجموعات
المتصفة بهذه الخاصية وحدها هي التي تنتمي إلى~$R$، ومن ثم لا يمكن أن
تكون~$R$ !!a{element} من~$R$، أي إن $R \notin R$. غير أن~$R$ لا يمكن
أن تكون وألا تكون في آن واحد !!a{element} من~$R$، فنقع في تناقض.

وبما أن افتراض $R \in R$ يؤدي إلى تناقض، يكون $R \notin R$. لكن هذا
يؤدي أيضًا إلى تناقض!{} فإذا كان $R \notin R$، كانت~$R$ نفسها متصفة
بخاصية~$R$ التعريفية؛ وعندئذ تكون~$R$ !!a{element} من~$R$ شأنها شأن
جميع المجموعات الأخرى غير المنتمية إلى نفسها. ومرة أخرى، لا يمكن أن
تكون في آن واحد !!a{element} من~$R$ وغير منتمية إليه.
\end{explain}
\end{tagblock}

\begin{digress}
كيف نؤسس نظريةً للمجموعات تتجنب الوقوع في مفارقة راسل، أي تتجنب
الادعاء \emph{غير المتسق} بأن المجموعة
$R = \Setabs{x}{x \notin x}$ موجودة؟ لا بد لنا من وضع بديهيات تحدد
شروطًا دقيقة جدًا للقول متى توجد المجموعات (ومتى لا توجد).

لا تفعل نظرية المجموعات الموجزة في هذا الفصل ذلك؛ فهي
\emph{ساذجة حقًا}. % حفظ أمر تشكيل غير منفذ من المصدر: \"
ولا تخبرنا إلا أن المجموعات تحقق الامتدادية، وأنه
إذا كانت لدينا بعض المجموعات، استطعنا تكوين اتحادها وتقاطعها، وما إلى
ذلك. ومن الممكن بناء نظرية المجموعات على نحو أشد صرامةً من هذا.
\oliflabeldef{cumul:::part}{سيُرجأ ذلك القدر من الصرامة إلى الجزء
\olref[cumul][][]{part}. أما الآن، فسنتابع على نحو ساذج، ونحاول بحذر
تجنب التناقضات. % حفظ أمر تشكيل غير منفذ من المصدر: \"
}{}
\end{digress}


\end{document}
